\section{Triển Khai}

Chương này trình bày cách BeVietnam được đưa lên môi trường chạy thật cho
giai đoạn pilot, từ cách bố trí dịch vụ trên máy chủ GPU đến cơ chế công
khai hệ thống ra Internet và cách bảo vệ các thông tin bí mật.

\subsection{Bố Trí Triển Khai}

Toàn bộ phần lõi phía server được gom về một máy chủ GPU duy nhất (1
card L40) để đơn giản hóa vận hành trong pilot. Trên máy này, một script
khởi tạo dựng lần lượt các dịch vụ: cơ sở dữ liệu, lưu trữ ảnh, mô hình
thị giác, AI Core và backend; cuối cùng là một đường hầm Cloudflare để
công khai backend ra một tên miền cố định. Web client chạy độc lập và
trỏ về tên miền này, còn mobile được phân phối dưới dạng APK cài trực
tiếp lên điện thoại. Hình~\ref{fig:deploy} mô tả bố trí đó.

\begin{figure}[H]
\centering
\resizebox{0.96\textwidth}{!}{%
\begin{tikzpicture}[
    >=Stealth,
    vm/.style={draw=primaryblue, thick, rounded corners=6pt, inner sep=8pt},
    svc/.style={draw=accentblue, fill=accentblue!8!white, rounded corners=4pt, minimum width=2.5cm, minimum height=0.85cm, align=center, font=\scriptsize},
    gpu/.style={draw=jade, fill=green!6!white, rounded corners=4pt, minimum width=2.5cm, minimum height=0.85cm, align=center, font=\scriptsize},
    edge/.style={draw=warnorange, fill=orange!8!white, rounded corners=4pt, minimum width=2.6cm, minimum height=0.85cm, align=center, font=\scriptsize},
    client/.style={draw=darkgray, fill=lightgray, rounded corners=4pt, minimum width=2.4cm, minimum height=0.85cm, align=center, font=\scriptsize}
]
% clients
\node[client] (web) at (0,1.0) {Web client};
\node[client] (mob) at (0,-1.0) {Mobile (APK)};
% edge
\node[edge] (cf) at (3.3,0) {Cloudflare\\Tunnel};
% VM box content
\node[svc] (be) at (7.0,1.6) {Backend\\FastAPI};
\node[svc] (ai) at (7.0,0.4) {AI Core};
\node[gpu] (vllm) at (7.0,-0.8) {vLLM\\Qwen2.5-VL};
\node[svc] (pg) at (10.3,1.6) {PostgreSQL};
\node[svc] (minio) at (10.3,0.4) {MinIO};
\node[gpu] (l40) at (10.3,-0.8) {GPU L40};

\begin{scope}[on background layer]
\node[vm, fit=(be)(vllm)(pg)(l40), label={[font=\small\bfseries, primaryblue]above:Máy chủ GPU (setup.sh)}] {};
\end{scope}

\draw[->, thick, darkgray] (web) -- (cf);
\draw[->, thick, darkgray] (mob) -- (cf);
\draw[->, thick, warnorange] (cf) -- node[above, font=\scriptsize]{https} (be);
\draw[->, thick, accentblue] (be) -- (ai);
\draw[->, thick, jade] (ai) -- (vllm);
\draw[->, thick, accentblue] (be) -- (pg);
\draw[->, thick, accentblue] (be) -- (minio);
\draw[->, thick, jade] (vllm) -- (l40);
\end{tikzpicture}}
\caption{Bố trí triển khai pilot. Toàn bộ dịch vụ server nằm trên một máy chủ GPU, được dựng tự động bởi script \texttt{setup.sh}. Cloudflare Tunnel công khai backend ra một tên miền cố định; cả web lẫn mobile đều truy cập qua đường hầm này nên hoạt động được trên mọi mạng.}
\label{fig:deploy}
\end{figure}

\subsection{Tự Động Hóa Khởi Tạo}

Việc đưa nhiều dịch vụ lên cùng một máy được gói trong một script khởi
tạo duy nhất, chạy tuần tự theo thứ tự phụ thuộc
(Bảng~\ref{tab:deploy_steps}). Cách này giúp một lệnh duy nhất có thể
dựng lại toàn bộ stack, đồng thời ghi cấu hình ra các file môi trường
riêng cho từng dịch vụ.

\begin{table}[H]
\centering
\caption{Trình tự khởi tạo của script triển khai}
\label{tab:deploy_steps}
\begin{tabularx}{\textwidth}{>{\raggedright\arraybackslash}p{0.9cm} >{\raggedright\arraybackslash}p{3.2cm} X}
\toprule
\textbf{Bước} & \textbf{Dịch vụ} & \textbf{Vai trò} \\
\midrule
1 & PostgreSQL & Khởi tạo cơ sở dữ liệu và chạy migration. \\
2 & MinIO & Dựng kho lưu trữ ảnh capture. \\
3 & vLLM & Tải và phục vụ mô hình thị giác Qwen2.5-VL trên GPU. \\
4 & AI Core & Khởi động dịch vụ agent (gồm Capture Judge). \\
5 & Backend & Khởi động API nghiệp vụ, kết nối DB, MinIO và AI Core. \\
6 & Cloudflare Tunnel & Công khai backend ra tên miền cố định. \\
\bottomrule
\end{tabularx}
\end{table}

Riêng vLLM~\cite{vllm_docs}, vì phải phục vụ mô hình thị giác, cần một
cấu hình đặc thù: chọn đúng model VLM, cho phép gửi tối đa hai ảnh trong
một yêu cầu (ảnh tham chiếu và ảnh người dùng), và đặt độ dài ngữ cảnh đủ
lớn vì token ảnh khá nặng. Một card L40 48\,GB đủ chỗ cho mô hình bản 7B
cùng KV-cache nên không cần thêm GPU.

\subsection{Bảo Mật và Cấu Hình Bí Mật}

Một nguyên tắc vận hành được tuân thủ chặt là \textbf{không đưa thông tin
bí mật vào mã nguồn quản lý phiên bản}. Các khóa dịch vụ (Foursquare,
OpenWeather, Goong, SerpAPI, token mô hình), bí mật của đường hầm và file
cấu hình thực thi đều được giữ ngoài Git. Backend chỉ lắng nghe nội bộ và
chỉ được công khai gián tiếp qua đường hầm Cloudflare, nên không mở cổng
trực tiếp ra Internet công cộng. Các ảnh tham chiếu tải về từ web phục vụ
việc xác minh cũng được lưu cache cục bộ và loại khỏi Git để tránh phát
tán nội dung có bản quyền.

\begin{note}
Việc gom toàn bộ server lên một máy GPU là lựa chọn có chủ đích cho
pilot: đơn giản, rẻ và đủ cho quy mô thử nghiệm. Khi mở rộng, các dịch vụ
trạng thái (PostgreSQL, MinIO) và dịch vụ tính toán nặng (vLLM) hoàn toàn
có thể tách ra các máy riêng mà không phải đổi kiến trúc, vì chúng đã
giao tiếp qua API và biến môi trường ngay từ đầu.
\end{note}
